% =====================================================================
\chapter{INTRODUCTION}
\label{ch:intro}
% =====================================================================

\section{Background}
\label{sec:background}

The last mile, the final leg from a pickup point to the customer, is widely
regarded as the most expensive and least efficient part of a parcel's
journey, and it has attracted a large body of operational research on new
delivery concepts \citep{boysen2021lastmile}. One of the most visible of
these concepts is \emph{crowdsourced delivery} (also called crowdshipping):
instead of relying only on employed drivers, a platform hands orders to
independent individuals who are paid per task. Reviews of the field document
both the rapid growth of such platforms and the new planning problems they
create, because the people who deliver are not employees and their capacity
and costs are not under the platform's control
\citep{alnaggar2021crowdsourced,savelsbergh2022challenges,savelsbergh2024update}.
Crowdshipping is also studied as a lever for more sustainable urban
logistics, since parcels can ride along on trips that take place anyway
\citep{mohri2023crowdshipping}.

In practice a delivery platform therefore combines three sources of
capacity \citep{archetti2016vrpod,luy2024workforce}, summarised in
Table~\ref{tab:classes}:
\begin{itemize}[leftmargin=*,itemsep=2pt]
\item a \emph{fixed-capacity fleet} (FD) operated under contract, which
serves any order at a known, published price;
\item \emph{gigworkers} (GW), who accept multi-stop routes through an
application. A gigworker starts from her current location, and her route is
\emph{open}: it ends at the last delivery;
\item \emph{occasional drivers} (OD), who carry parcels along a personal
trip they would make anyway. Their route must end at their own destination,
and the extra time they are willing to spend is limited by a \emph{detour
budget}.
\end{itemize}

\begin{table}[t]
\centering
\caption{The three sources of last-mile capacity considered in this thesis}
\label{tab:classes}
\small
\begin{tabular}{P{2.5cm}P{4.3cm}P{2.0cm}P{5.0cm}}
\toprule
Class & Route geometry & Strategic? & Private information\\
\midrule
Gigworker (GW) & Open route from the current location, ends at the last delivery & Yes & Value of time $\theta_i$ (one scalar)\\
Occasional driver (OD) & From origin to personal destination, limited by a detour budget & Yes & Value of time $\theta_i$ (one scalar)\\
Fixed fleet (FD) & Not routed in the model; serves any order & No & None: public price $q_o$ per order\\
\bottomrule
\end{tabular}
\par\smallskip\footnotesize Source: taxonomy adopted from \citet{luy2024workforce}.
\end{table}

The fixed fleet has a known price. The other two sources do not. How much
an hour of driving is worth to a particular driver --- her \emph{value of
time} --- is private information, and it determines her cost of serving a
route. A platform that pays drivers what they claim invites them to
overstate it. The classical remedy from mechanism design is the
Vickrey--Clarke--Groves (VCG) mechanism: the platform selects the
allocation that minimises total reported cost and pays each winner her
externality on the others. Under this payment rule, reporting the true cost
is a dominant strategy, that is, the best action for a driver regardless of
what the others report \citep{nisan2007agt}. In routing applications the
mechanism is implemented in four steps \citep{zou2022mechanisms,li2026auction}:
the platform generates candidate routes, collects one scalar bid per driver,
solves a set-partitioning winner-determination problem (WDP), and re-solves
it once per winner to compute the Clarke-pivot payments.

This thesis studies such a procurement mechanism for a platform that buys
from \emph{both} crowd classes at once, with the fixed fleet as a publicly
priced outside option. Studying the two crowd classes together matters because their capacities
are complementary: an occasional driver is cheap but only for orders that
lie along her trip, whereas a gigworker can go anywhere but charges for all
of her working time.

\section{Problem Statement}
\label{sec:problem}

\subsection{The need for a bid-independent route pool}

For the VCG mechanism above to remain truthful, the set of routes among
which the platform chooses must be fixed \emph{before} any bid is seen. The
mechanism is then VCG on a bid-independent range, and dominant-strategy
incentive compatibility follows from the maximal-in-range argument of
\citet{nisan2007computationally}. If the platform let the bids influence
which routes are generated --- for example by generating routes with column
generation, whose pricing step depends on dual values and hence on the bids
--- a driver could change the set of available routes by changing her bid,
and truthfulness would no longer be guaranteed.

\subsection{The pool is the bottleneck, and the two classes differ}

The price of this requirement is the size of the route pool. For an
occasional driver, the detour budget anchors the search: adding orders to a
bundle can only lengthen the detour, so a bundle whose detour already
exceeds the budget can be discarded together with all its supersets before
any visiting sequence is enumerated \citep{li2026auction}. A gigworker has no
such anchor. Her route is open and is limited only by time windows. Two
orders may each be servable on their own while the pair is not, and a bundle
of distant orders remains feasible whenever the windows are wide. Her pool
therefore grows as $\Theta(n^B)$ in the number of orders $n$ for a bundle cap
$B$. On the experimental instances of this thesis, the pool of an occasional
driver is only 0.7--1.7\% of that of a gigworker, and on five instances only
0.4--1.5\% of all generated routes are optimal for \emph{any} of 1,000 sampled
bid profiles. Most of the pool is never used, yet the platform cannot tell
in advance which part.

Removing routes is not straightforward either. Every cheap reduction rule
that was tried during the development of this thesis either deleted routes
that some bid profile needs, or pruned almost nothing
(Section~\ref{sec:selection} lists them). Two obstacles explain why:
\begin{enumerate}[leftmargin=*,itemsep=2pt]
\item any rule that reads the bids changes the allocation range and can
destroy truthfulness;
\item any rule that reasons about \emph{other} drivers must, in the worst
case, solve the winner-determination problem itself, which is NP-hard
(Remark~\ref{rem:scope}).
\end{enumerate}
What remains are \emph{local, bid-independent} rules, which look only at
the pruned driver's own data and the public prices. No study in the
literature characterises how far such rules can go
(Chapter~\ref{ch:related}).

\subsection{Research questions}

The problem leads to one central question and four technical
sub-questions. This proposal answers TQ1 in Chapter~\ref{ch:solution};
TQ2--TQ4 are planned for the final thesis (Section~\ref{sec:sd-planned}).

\begin{quote}
\emph{How far can a route pool be pruned, before bids are observed, by a rule
that uses only the pruned driver's own data, without changing the outcome of
the VCG mechanism?}
\end{quote}

\begin{description}[leftmargin=1.2cm,style=nextline,itemsep=2pt]
\item[TQ1] Which routes can a local bid-independent rule safely delete, and
is there a smallest pool that every such rule must keep?
\item[TQ2] How large can that smallest pool be, and does it contain routes
that are never useful?
\item[TQ3] Can the pruning be moved inside the route enumeration so that
fewer partial routes are explored, without losing safety?
\item[TQ4] Can the repeated counterfactual solves required by the payments be
decomposed into smaller problems, and when does this fail?
\end{description}

Since the mechanism is meant to be used by a platform, the thesis also asks
five empirical questions about its economic and computational behaviour,
which are answered by the computational study (Chapter~5):

\begin{description}[leftmargin=1.2cm,style=nextline,itemsep=2pt]
\item[RQ1] Does letting gigworkers and occasional drivers bid jointly lower
system cost compared with using either class alone or consulting them in
sequence?
\item[RQ2] Does endogenous bundling --- letting the market choose among
platform-generated bundles --- add value compared with single-order routes
and with a fixed bundling heuristic?
\item[RQ3] What does truthful procurement cost compared with a
full-information benchmark, pay-as-bid and posted prices?
\item[RQ4] How much of the pool does the frontier remove, does it change any
payment, and what does it save in computation?
\item[RQ5] Are the conclusions robust to time windows, detour budgets,
fixed-fleet prices, spatial clustering, type dispersion and instance size?
\end{description}

\section{Objectives of the Study}
\label{sec:objectives}

\subsection{General and specific objectives}

The general objective is to design, analyse and evaluate a truthful
procurement mechanism in which gigworkers and occasional drivers bid jointly
for platform-generated route bundles, with the fixed fleet as a publicly
priced outside option, and to characterise exactly how far its route pool
can be reduced without affecting the mechanism. The specific objectives are:
\begin{enumerate}[label=O\arabic*.,leftmargin=1.2cm,itemsep=2pt]
\item To formulate the procurement problem of a static dispatch session as a
mechanism: a cost model for the two crowd classes, a winner-determination
model, and Clarke-pivot payments.
\item To develop a complete, bid-independent route-generation algorithm for
both driver classes (Algorithm~A).
\item To define a local pruning rule (the local frontier), prove that it is
safe (TQ1), and determine whether it is minimal and how large it can be
(TQ2).
\item To develop and prove a label rule that applies the frontier reasoning
during route generation (TQ3), and to establish when payment counterfactuals
decompose (TQ4).
\item To implement the complete pipeline and verify every component against
an independent brute-force oracle.
\item To evaluate the mechanism in a pre-registered computational study
(RQ1--RQ5).
\end{enumerate}

\subsection{Beneficiaries and expected outputs}

The results address three groups.
\emph{Delivery platforms} obtain a procurement procedure that remains
correct when the crowd becomes competitive, and a certificate, computed
before any bid is collected, that tells them whether an auction is worth
opening and which drivers need to be invited. \emph{Drivers} benefit
because truthful reporting is their best strategy, so they need no
strategic reasoning, and because a winner is never paid less than her
reported cost. \emph{Researchers} in routing and mechanism design obtain a
provably safe way to reduce route pools before bidding, which applies to any
set-partitioning auction with a priced outside option.

The expected outputs are:
\begin{itemize}[leftmargin=*,itemsep=2pt]
\item a three-stage procurement pipeline (route generation, winner
determination, payments) with a provably safe pruning step in between;
\item a theorem, with full proof, that pruning the pool to the local frontier
is safe, followed in the final thesis by results on its minimality and
worst-case size, the label rule and payment decomposition;
\item a tested implementation in Python with the CPLEX solver, a synthetic
instance generator and locked, hashed experiment parameters;
\item a computational study that quantifies when joint bidding, bundling and
truthful payments matter.
\end{itemize}

\subsection{Expected scientific contributions}

\begin{enumerate}[leftmargin=*,itemsep=3pt]
\item \emph{The local pruning frontier} (Theorem~\ref{thm:frontier}, proved
in this proposal). An envelope compares each route with the cheapest subset
route of the same driver whose unserved orders are sent to the fixed fleet,
uniformly over the admissible report range. Deleting every route outside the
resulting frontier preserves the optimal value of the winner-determination
problem, every counterfactual value and therefore all VCG payments, because
every deleted route has a substitute that is itself retained.
\item \emph{Minimality and the price of locality} (planned). Every route in
the frontier should be indispensable in some instance that looks the same
from the driver's own data, and a single gigworker's frontier can contain
$\Theta(n^B)$ routes that are never optimal.
\item \emph{A safe label rule} (planned). Rollback pruning
\citep{lozano2016exact} cannot be transferred directly, because working time
is priced and waiting is allowed; a correction for waiting makes it safe.
\item \emph{Payment decomposition with a matching negative result}
(planned). With a separable outside option the counterfactual solves
decompose along the components of a conflict graph; a shared capacity on the
outside option breaks the decomposition.
\item \emph{A pre-registered computational study} of RQ1--RQ5, in which all
parameters were locked and hashed before the main runs.
\end{enumerate}

Two things are \emph{not} claimed. The incentive property is not new: the
mechanism is exact VCG on a bid-independent range, and its truthfulness is
classical. Nor is the three-class driver taxonomy, which is adopted from
\citet{luy2024workforce}.

\subsection{Design requirements}

The mechanism and its implementation must satisfy the requirements in
Table~\ref{tab:requirements}. Requirements R1--R7 are correctness
requirements that the design must meet exactly; R8 and R9 are practical
requirements.

\begin{table}[t]
\centering
\caption{Design requirements of the procurement system}
\label{tab:requirements}
\small
\begin{tabular}{P{0.8cm}P{3.6cm}P{9.6cm}}
\toprule
No. & Requirement & Description\\
\midrule
R1 & Truthfulness (DSIC) & Reporting the true value of time is a dominant strategy for every driver of both classes.\\
R2 & Individual rationality & A winner is paid at least her reported cost; a loser pays and receives nothing.\\
R3 & Service coverage & Every order is served exactly once, by one selected route or by the fixed fleet.\\
R4 & Exactness & Every winner-determination and counterfactual problem is solved to proven optimality (zero gap).\\
R5 & Bid independence & No step before report collection reads a report; the pool is fixed and hashed before bidding.\\
R6 & Deterministic tie-breaking & Ties among optimal solutions are broken by a rule that does not depend on reports.\\
R7 & Safe pruning & Any pool reduction leaves the optimal value, every counterfactual value and every payment unchanged for every admissible report profile.\\
R8 & Tractability & The complete pipeline finishes within 300 seconds per instance up to 30 orders on a single CPU thread, and route generation remains tractable for instances of 100 orders or more.\\
R9 & Reproducibility & Parameters are fixed and hashed before the main runs; every component is checked against an independent oracle.\\
\bottomrule
\end{tabular}
\end{table}

\section{Scope and Limitation}
\label{sec:scope}

\subsection{Scope}

\paragraph{Problem setting.} The study considers a single, static and
deterministic dispatch session: all orders and drivers are known when the
session opens, travel times are deterministic, and the mechanism is run once.
Every order must be served (mandatory fulfilment), each driver receives at
most one route (unit demand), and routes carry at most $B$ orders. Each
driver has one private parameter, her value of time; all other driver data
are public or committed before the session opens.

\paragraph{Instances and data.} All experimental instances are synthetic.
They are produced by an instance generator on a plane with dispersed pickups
and deliveries, time windows of 120 minutes, a detour budget of 30 minutes
for occasional drivers, 5 minutes of service per stop and a speed of
20\,km/h. The distance coefficient is $\kappa=1$ cost unit per kilometre,
values of time are drawn from $U[18,25]$ cost units per hour for both
classes, and the fixed-fleet price of order $o$ is
$q_o=8+3\max\{0,\mathrm{dist}_o-1\}$, where $\mathrm{dist}_o$ is its
pickup--delivery distance in kilometres. The main grid crosses five levels of
an \emph{alignment} parameter (how strongly occasional drivers' trips overlap
the demand region), three instance sizes $n\in\{10,15,20\}$ and four supply
mixes with two or three drivers of each class, with 25 replications per
cell (1,500 instances). Robustness tests extend the size to $n=30$. The
bundle cap is $B=3$, with $B=4$ as a sensitivity case for $n\le 15$, and the
report domain is $\Theta=[18,25]$. No time frame of data collection applies,
since no field data are used.

\paragraph{Instance size.} Small instances are not the limit of the method.
A separate scaling study has already run the route-generation algorithm
(Algorithm~A) on instances of the same generator with up to $n=100$ orders,
bundle caps $B\in\{2,3,4\}$ and time windows of 30 to 240 minutes
(Section~\ref{sec:sd-algA}). The results chapter will extend the evaluation of
the complete mechanism to large instances with $n\in\{50,75,100\}$ orders,
and to larger instances where the time budget allows, so that its
behaviour at the scale of a real dispatch session can be assessed.

\paragraph{Tools.} The pipeline is implemented in Python~3.7.7. The
winner-determination and counterfactual problems are solved with
IBM ILOG CPLEX~12.10 on one thread with zero relative and absolute gap
tolerances.

\subsection{Main assumptions}

\begin{enumerate}[label=(\roman*),leftmargin=1cm,itemsep=2pt]
\item Static, deterministic, single-session procurement.
\item Mandatory fulfilment: each order is served exactly once, by a route or
by the fixed fleet.
\item The fixed fleet is uncapacitated and serves each order at a public,
fixed price $q_o$.
\item Unit demand: each driver receives at most one route.
\item The allocation range (bundle cap $B$, route generation and any
pruning) is fixed before bids.
\item Single-parameter private information: a driver's cost is affine in
her value of time, $C_{ir}=K_{ir}+\theta_iW_{ir}$
(Section~\ref{sec:sd-costs} explains why this is an empirical approximation
rather than a mere convenience).
\item Occasional drivers' destinations and detour budgets are public, or are
committed before the session opens.
\item Travel distances and times satisfy the triangle inequality; arriving
before a time window opens is met by waiting.
\end{enumerate}

\subsection{Limitations}

\begin{itemize}[leftmargin=*,itemsep=2pt]
\item \emph{Synthetic data.} All instances are synthetic and static; no
field data from a delivery platform were available. The fixed-fleet price
was calibrated on a pilot without corridor concentration, and that
calibration did not carry over to the main grid, so the alignment levels do
not order the regimes monotonically by how competitive the crowd is. Results
are therefore reported by regime and observed fixed-fleet share.
\item \emph{Information structure.} Treating occasional drivers' destination
and detour budget as public is a restriction: \citet{li2026auction} let
carriers report four parameters. Private costs are single-parameter and
affine in the report.
\item \emph{Report domain.} Pruning requires a report domain announced in
advance. The pay-as-bid benchmark inherits that bound, which favours it.
\item \emph{Scale.} The economic experiments completed so far cover
$n\le 30$; at $n=100$ only route generation has been measured so far. The
largest bundle cap ($B=4$) was not evaluated at $n=20$, and the scaling study
shows that it becomes intractable before $n=50$, so large instances are
limited to $B\le 3$. The robustness experiment is centred on the least
competitive regime.
\item \emph{Implementation.} The label rule has been integrated only into an
experimental build of the route generator, and in pure Python its extra
bookkeeping currently makes route generation slower
(Section~\ref{sec:sd-label}). The payment decomposition requires a
separable outside option.
\item \emph{Out of scope.} The study makes no claims about online arrivals,
stochastic travel times, budget balance, collusion, false-name bids or
fairness.
\end{itemize}

\section{Structure of the Thesis}

Chapter~\ref{ch:related} introduces the theoretical foundations and reviews
the literature on crowdshipping, procurement mechanisms, column pool
reduction, label dominance and VCG payment computation, and describes the
candidate solution methods. Chapter~\ref{ch:method} compares these
candidates, selects the approach and presents the design of the proposed
system. Chapter~\ref{ch:solution} develops the solution: the mathematical
model, the route-generation algorithm, the local pruning frontier and its
safety theorem, payment computation, the results planned for the final
thesis, a complete numerical example, the implementation and a comparison of
alternatives. The computational study (Chapter~5) and the conclusions
(Chapter~6) follow. The proof of the key lemma is in
Appendix~\ref{app:proofs}.
